% !TeX program = xelatex
\documentclass[11pt]{article}
\usepackage[a4paper,margin=2.3cm]{geometry}
\usepackage{amsmath,amssymb,booktabs,tabularx,array}
\usepackage{fontspec}
\usepackage[unicode,hidelinks]{hyperref}
\defaultfontfeatures{Ligatures=TeX}
\setmainfont[AutoFakeSlant=0.2]{Noto Serif CJK SC}
\setsansfont{Noto Sans CJK SC}
\setmonofont{Noto Sans Mono CJK SC}
\XeTeXlinebreaklocale "zh"
\XeTeXlinebreakskip=0pt plus 1pt minus 0.1pt
\setlength{\parindent}{2em}
\setlength{\parskip}{0.35em}
\setlength{\emergencystretch}{3em}
\renewcommand{\arraystretch}{1.2}
\setcounter{tocdepth}{1}
\newcommand{\term}[1]{\textbf{#1}}
\newcommand{\source}[1]{\par\noindent{\footnotesize 详见：\path{#1}}\par}
\title{Model 0826 Plus：怎样拟合，怎样检验\\\large 从个体参数搜索到可信的行为与规则解释}
\author{内部工作说明}
\date{2026年9月18日修订}
\begin{document}
\maketitle
\noindent\term{先说结论。}正式拟合适合采用“每个人先做一轮有覆盖的搜索，再按问题追加计算”的方式。
有的人需要多找几组参数，有的人需要把已有参数的得分算得更准。
参数选定以后，再单独提高最终预测和规则信念的精度。
我们已用少量被试做了初步检验，还没有把它作为全体被试的正式默认方案。

本文是 \texttt{model\_0826.tex} 的实践说明。主文档负责模型方程；本文回答计算怎么安排、结果能解释到什么程度、后续证据怎样组织。
本次精简原有章节，保留行为、恢复、口述和自主生成的检验要求。
旧稿及其 PDF 已随本轮结果归档，历史计划不再作为当前运行状态。

\section{先分清三件事：找参数、算得分、看状态}
可以把参数拟合理解为寻找一处低谷。\term{搜索覆盖}决定我们走过多少条路；\term{评分精度}决定能不能分清两处位置的高低；\term{状态精度}决定最后画出的内部过程是否稳定。
三者有关，但不能用同一种追加计算解决。

\begin{center}\small
\begin{tabularx}{\linewidth}{p{2.1cm} X X}
\toprule
问题 & 通俗含义 & 应增加的计算\\
\midrule
搜索覆盖不足 & 分数算得比较清楚，但还没找到已知更好的区域 & 真正不同的起点、不同方向的参数调整\\
评分波动过大 & 同一组参数重算几次，候选优劣仍分不清 & 检查波动来源，针对候选校准粒子和重复数\\
状态输出不稳 & 平均分数已较稳定，规则信念曲线仍随随机运行改变 & 固定参数后增加输出精度，并检查状态可识别性\\
\bottomrule
\end{tabularx}
\end{center}

粒子滤波（particle filter，PF）用有限样本近似模型内不可见的随机过程。
粒子数 $R$ 是一次计算保留多少条可能过程；重复数 $B$ 是用多少组独立随机种子重算。
它们都是\term{计算预算}，不是人的工作空间容量、记忆能力或学习次数。
$R$ 和 $B$ 也不能只按乘积互换：很多次低粒子运行，并不自动等价于一次高粒子运行。

前两轮试点给出了两个具体例子。选择得分使用负对数似然（NLL），越低表示拟合越好。
S129 的一批新参数比历史好参数的平均 NLL 高约 0.013，差距在数值复核中仍清楚可见。
这里需要先改善“找在哪里”。S229 的某些固定参数即使重复计算，得分和规则信念仍有较大波动。
这里需要先检查“算得有多准”。这些发现来自个案，不能直接变成整个 condition 的统一预算。

\section{正式拟合建议怎样执行}
\subsection{第一步：所有被试先获得共同的最低搜索覆盖}
使用同一套模型、允许搜索的参数范围、评分定义和输入检查，先安排若干真正不同的起点。
“两个起点”不能只是工作空间开关不同，而其他参数完全相同，第一轮扫描后又合并成同一路线。
起点和后续调整应覆盖记忆、搜索控制、规则精度与执行结构等不同方向。

先用适中的计算精度发现候选，同时保留几组不同的较好参数。
不要求把粗网格、细网格都完整扫描；也不只沿着当前最低分的一个点反复微调。
低预算有用的前提是，它没有经常把真正有竞争力的区域过早淘汰。
这需要用高一些的预算抽查，不能因为运行很快就假定它足够。

\subsection{第二步：只对少数有竞争力的候选认真比较}
搜索结束或准备停止时，选出少数候选，用独立于搜索的种子比较。
候选之间尽量使用相同种子作配对比较，以减少无关的随机差异。
若要检验新搜索是否接近一个历史方案，先在筛选阶段各自选定代表，再用另一组确认种子比较。
不能看完确认结果后重新挑选代表，再把同一分数当作独立确认。

复核的目的，是防止保留一个在搜索中偶然获得低分的候选。
它不要求每个候选都精算，也不要求分数接近时强行选出唯一赢家。
几组较好参数的得分在当前精度下无法区分，就一并保留，称为近优候选；再检查它们是否给出不同的心理解释。

\subsection{第三步：按诊断追加计算，而非人人重复同一套流程}
\begin{center}\small
\begin{tabularx}{\linewidth}{X X}
\toprule
发现什么 & 接下来怎么做\\
\midrule
得分比较可靠，换方向仍找到明显更好的候选 & 继续拓展搜索；单纯增加粒子不能替代找参数\\
候选顺序随重复改变，差异与数值误差相当 & 先校准评分精度；暂不把随机波动当优化进步\\
搜索已平稳，但追加的一批独立方向候选仍能明显改善 & 原来的停止信号不可靠，继续搜索并修订停止规则\\
搜索平稳、评分清楚，额外挑战也未显示超出容差的改进 & 可在已检验的范围内暂时停止参数搜索\\
预算已用完，但上述检查没有通过 & 标记为尚未充分解决；记录缺口，不能写成已收敛\\
\bottomrule
\end{tabularx}
\end{center}

\term{搜索平稳只是必要证据之一。}如果两轮都没有改进，但新方向很容易找到更好的参数，前两轮可能只是在一个局部区域内打转。
因此，停止规则要经受额外搜索的挑战。
若改善的数值区间仍很宽，“没有显著改善”也不能证明“改善已小到可以忽略”。

有历史好参数的被试可以从那里继续搜索，节省重复发现的成本。
但报告必须区分“通用起点找到的结果”和“借助历史参数继续改进的结果”。
后者不能用来证明同样的通用策略对新被试也足够。

公平的要求是\term{相同质量标准}，而不是每个人消耗相同 CPU 时间。
以后比较多个模型时，每个模型也应达到对等的搜索与评分质量。
不能因某个简化模型分数暂时较差，就提前停止它的搜索。

\subsection{第四步：参数选定后，再提高最终输出精度}
用于找参数的预算与用于最终状态图的预算可以不同。
最终预测、规则信念、活跃概率等，应在固定参数下用足够精度重新生成，并记录对应种子和参数来源。
如果多个近优参数的状态解释不同，还需要报告参数不确定性；增加同一个参数点的粒子，无法消除这种差异。

这套安排保留必要的复核，但把昂贵计算集中在真正有竞争力的参数和最终需要的输出上。
它并未承诺“所有人都只用少量粒子即可”，也没有把粗搜、细搜、复评固定为必须逐人完整执行的三大阶段。

\subsection{怎样理解目前使用的容差}
试点暂以平均 NLL 差 0.005、逐试次选择概率的均方根误差（RMSE）0.02 等作为计算诊断门槛。
它们是探索性的实用容差，不是普适统计标准，也不是正式置信集。
对于 S229 的 1087 个评分试次，平均 NLL 差 0.005 相当于总 NLL 差约 5.44；
如果后续模型比较关心更小的总分差，评分精度必须相应收紧。

数值区间目前通过重采样 PF 种子得到，只描述固定数据、固定候选、有限粒子下的计算不确定性。
它不包括参数搜索遗漏、被试抽样或模型错误，也不能证明已经找到全局最优。
一个 $B=16$ 结果分成两半，是 8 对 8 的检查，不等于两次各含16次独立重复的复核。

\section{目前做到了哪里，本轮还在检验什么}
\subsection{任务与版本边界}
\begin{center}\small
\begin{tabularx}{\linewidth}{l l X X}
\toprule
论文任务 & condition & 反馈与实现 & 当前解释边界\\
\midrule
Task 1 & 1 & 二分类、0/1反馈、29条规则 & 已有代表性拟合和预算试点；全体正式流程待冻结\\
Task 2 & 3 & 四分类、0/0.5/1反馈、116条规则及配对学习 & PMH已接入；本轮检查完整序列搜索，恢复与标签支持审计待完成\\
Task 3 & 2 & 四分类、0/1反馈、116条规则 & 已有S229拟合及数值诊断；预算不能直接照搬二分类\\
\bottomrule
\end{tabularx}
\end{center}

Task 编号和 condition 编号不同。当前 Fig2 的 b/c/d 分别对应 condition 1/3/2；S229 对应 d。
Condition 2 的错误反馈只排除所选类别，不能推定另一个唯一正确类别。
其规则精度更新保留既有二值反馈形式，这是明确的模型假设，仍需科学验证。
Condition 3 则显式学习未知反应配对，处理部分正确反馈，不能用 condition 2 的二值反馈入口替代。

这次工作调整搜索和数值诊断，不修改任何 condition 的认知方程。
等价缓存、似然合并及并行属于工程加速；改变候选、粒子或种子属于估计方案调整，可能改变最终拟合参数与状态。
不能把这两类操作都笼统称为结果完全相同的加速。

\subsection{已有两轮试点说明了什么}
前两轮把稠密搜索压缩为有上限的低预算搜索，并尝试更分散的起点和联合调整。
S129 的计算更快，但仍落后于已知历史候选，且到达上限时还在改善。
S229 固定四组参数校准到 $R=128,B=32$，仍未全面通过评分与状态门槛。
这说明简单削减预算尚不能直接替代历史拟合；具体耗时与候选记录保留在各轮报告中。

\subsection{本轮检查的三个问题}
本轮将逐个参数方向的长距离调整与邻近调整结合，减少在很大的联合参数块内盲目抽样。
S129 分别检查通用起点和历史好参数续搜；S229 检查高波动候选、时间段、有效粒子数和重采样，并作较高粒子数对照。
另外选择 S103、S203、S301，均按所在 condition 的试次数接近中位数确定，同距离取较小编号，不根据模型分数挑选。

每名搜索被试都接受一批额外候选的挑战。
若中途达到“两轮改善很小”，先冻结当时的候选，后面的搜索仍继续作为停止规则的检查。
主比较在独立筛选后固定，随后才进行确认。
本轮结果与后续决定见配套结果报告；少量个案足以揭示失败方式，不能估计全体被试的通过比例。
\source{results/model_0826/search_stopping_pilot_20260918/README.md}

\term{S129 的进展。}通用搜索加追加挑战得到的候选，独立确认 NLL 为0.5075，历史参照为0.5074。
补充配对分析的差值区间约为 $[-0.0030,0.0029]$，在本轮探索容差内接近。
但两者的规则信念与策略仍不同；从历史参数继续搜索的候选虽有更低点估计，也尚未确认优于历史参照。
现在可以说搜索找到了更有竞争力的区域，不能说已确定唯一参数或内部过程。

\term{S229 的波动。}四个候选的旧种子结果均逐元素完全复现。
高波动候选的平均有效粒子数并不特别低，却有更多试次出现粒子权重强烈集中；误差还会跨试次相关。
对两个高波动候选，把粒子数从128增到256、重复数固定16，评分标准误分别下降约40\%和53\%，规则信念仍未稳定。
因此应针对候选和输出目标安排计算，不能只看平均有效粒子数，也不能把评分改善当作状态已经可靠。

\term{新增被试与停止检查。}S103 的评分比较通过本轮门槛，但搜索尚未形成平台；S203 和 S301 的评分精度未通过，应先校准关键候选的得分。
四名搜索被试均未在预设轮数内触发两轮平台，因此本轮尚不能估计实际早停的误停率。
这套规则仍需在确实达到平台的个案上接受后续搜索挑战，不能现在就作为已经验证的正式默认流程。
Condition 3 也尚未完成全部规则与配对联合状态的精度和恢复检查。

\section{读懂拟合结果：得分好不等于内部过程已被证明}
\subsection{我们到底在给什么评分}
对同一参数点，先把 $B$ 个 PF 重复的逐试次概率平均，再取对数：
\begin{equation}
 \bar p_t(c)=\frac1B\sum_{b=1}^B\widehat p_t^{(b)}(c),\qquad
 L(\theta)=-\frac1{|\mathcal T|}\sum_{t\in\mathcal T}\log\bar p_t(y_t).
\end{equation}
NLL 越低，模型给实际选择的概率总体越高。
这里不是先分别求每次 NLL 再平均，也不是选择最贴近数据的一条潜在路径。

当前约定保留完整 trial 序列递推，但首试次不进入默认评分掩码。
因此 S129 的256个试次中有255个计分，S229 的1088个中有1087个计分。
原始时间顺序、session/block边界和反应编码保持不变；不能为了稳定分数删掉高波动试次。

每次选择必须由作答前状态预测，再用已观察的选择与反馈更新状态。
即使每一步递推严格遵守时间顺序，\term{只要参数用整段数据拟合，成绩仍然是样本内拟合}。
独立 PF 种子只使计算复核独立，不会让同一份行为数据变成留出测试集。

\subsection{规则信念、活跃概率与完整路径不要混用}
规则信念表示各规则获得多少支持，归一化后总和为1；活跃概率表示规则进入工作空间的可能性，其总和通常等于容量。
两者不能相互替代。持续执行结构 $\chi=1$ 下还有执行规则；$\chi=0$ 的信念混合读出没有这个变量。
不能把最大信念规则的变化直接称为被试的执行切换。

作答前状态只使用此前信息；过滤后状态吸收了当前选择，需注明是否也已吸收反馈；
利用整段记录回溯的祖先路径还使用了后续信息。
这些输出回答不同问题，图和分数要使用一致的信息时点。
最佳完整路径适合辅助解释，不能每个试次挑一条不同的路径拼成“真实认知过程”。

较稳定的平均 NLL，也可能对应不同的规则信念或内部策略。
所以结果应分开回答：选择预测是否稳定、参数是否近似等价、状态解释是否一致。
真实被试没有已知的心理状态真值，历史拟合也只能作为参照，不能当作真值。

\section{拟合完成后，还需要哪些证据}
\subsection{行为：分开检查拟合、留出预测与自主生成}
样本内检查用于发现问题：学习曲线、类别偏向、错误模式、逐类别概率和极低概率选择。
正式泛化分析需按时间前缀拟合，在共同后缀评价冻结参数；比例应提前确定，不能随机打乱学习 trial。
若预测后继续吸收已发生的后缀选择与反馈，应称为固定参数的在线一步预测。

主比较以同一被试、同一评价试次的 NLL 差为单位，同时使用均匀、类别偏向及适当的简单学习基线。
不同类别数下先做任务内比较，不直接把二分类和四分类的原始 NLL 排名解释为模型优劣。
群体结果每个点代表一名被试；PF种子、重叠滑窗和时间段不能冒充独立被试。
失败或未稳定的个案也需要报告，不能只展示成功例。

\subsection{恢复：能找回什么，就解释到什么程度}
用模型自己生成的选择与反馈建立已知真值，再重新拟合。
分别检查模块架构、执行结构、连续参数和潜在状态；其中一项通过不代表其余都通过。
生成轨迹不能先平均成一个虚拟被试，每条轨迹应独立拟合。

参数恢复还要检查边界堆积、零值与正值混淆，以及参数之间的补偿。
状态恢复需在相同信息时点比较规则信念、活跃集合、执行与切换。
若多种状态产生近似相同选择，应报告不可区分范围，不把其中一条当作唯一心理解释。
现有恢复设计和已生成数据不等于完整恢复结论已经成立。

\subsection{口述：检验它是否增加选择之外的信息}
口述没有进入当前选择目标函数，但规则目录和编码开发曾参考口述，因此应保留开发与验证来源记录。
正式分析冻结编码尺度和规则目录，明确有效新报告、缺失、不可解析与沿用历史报告之间的区别。
当前报告在作答后、反馈前，模型状态的对齐时点也应事先确定。

可以比较完整规则目录上的模型分布 $q_t(h)$ 与口述分布 $o_t(h)$：
\begin{equation}
 O_t=\sum_h\min\{q_t(h),o_t(h)\}=1-\tfrac12\sum_h|q_t(h)-o_t(h)|.
\end{equation}
重叠接近1说明两种分布接近，不等于规则识别准确率，也未自动扣除机会一致性。
仅仅提到同一个特征，或两者都不支持目标规则，也不能证明双方采用同一条规则。

需要与静态偏好、共同刺激和选择所能解释的对应关系，以及保留时间结构的零模型比较。
主评分使用固定完整目录；仅在模型活跃集合内重归一化口述会改变比较问题。
在选择拟合近似相同的候选之间，口述能否进一步区分，是更直接的附加约束检验。
若用某批口述挑选模型，就需要新的留出报告或将结论明确标为探索性。

\subsection{自主生成与模型比较：让解释接受更强的挑战}
自主生成需要模型自己选择，环境据此给反馈，再由模型自己的历史更新。
继续输入真实人的选择与反馈是条件化重放，不是自主生成。
比较多条完整生成轨迹中的学习平台、改善、反转、持续错误和达到学习标准的时间分布，不能只挑最像人的一条。
有限日程结束仍未达标时，要记录删失，不能当作真实完成时间。

P、PM、PH、PMH 比较用于检验记忆与工作空间搜索的贡献；P 仍包含信念学习和动态规则精度，不能称为没有学习。
四种结构都保留动态精度，因此要证明动态精度本身有贡献，还需专门对照。
各模型应独立重拟合，采用共同评价集和对等的计算质量。
固定完整模型参数后关掉一个模块，只是敏感性检查，不能代表简化模型能达到的最佳表现。

模型内“错误累积后搜索增多”部分由方程决定。
要解释人的机制，还需独立口述变化、未来行为或其他测量的约束。
拟合、消融和事件相关本身不构成人类心理过程的因果证明。

\section{论文里怎样呈现，哪些内容留在补充材料}
正文的阅读顺序可以保持简单：先展示模型如何形成选择，再展示三个任务的个体过程，随后回答能否推广、口述是否增加约束、哪些结构有贡献。
下面是 Fig2 的工作分工，不表示各 panel 的正式数据已齐备。

\begin{center}\small
\begin{tabularx}{\linewidth}{l X X}
\toprule
Panel & 回答的问题 & 必须保留的说明\\
\midrule
a & 一次试次如何运行 & 认知更新与PF推断分开；作答与反馈时序清楚\\
b/c/d & 三任务个体过程 & condition 1/3/2；完整trial轴；预测、信念与口述时点一致\\
e & 个体结果能否推广 & 全体被试的共同评价集，每点一人，报告未稳定与缺失\\
f & 口述增加了什么约束 & 完整目录、共同输入对照与时间结构零模型\\
g & 哪些结构有贡献 & 少数核心模型的公平重拟合比较\\
\bottomrule
\end{tabularx}
\end{center}

补充材料承担覆盖和稳健性：全部个体图、分类校准与残差、预算和停止检查、结构与参数恢复、状态恢复、口述编码敏感性、扩展消融、自主生成分布。
决定主论点的关键证据仍要在正文给出摘要。
三任务的目录和试次数不同，不能用相同规则编号或拉伸后的trial轴制造可比性。
当前布局与其他主图的口述分析应避免重复，最终编号随全文安排确定。

\section{后续执行顺序与复现记录}
\begin{enumerate}
\item \term{先补足本轮发现的缺口。}对关键近优候选校准评分精度；对评分已清楚的个案继续有限搜索，并在真正出现平台后检验是否会误停。
\item \term{据此冻结正式拟合方案。}确定最低覆盖、追加计算条件、容差、计算上限和未解决个案的处理办法；按条件检查恢复与反馈语义，不把一个条件的成功直接推广到另一个条件。
\item \term{再做正式拟合与相应验证。}按冻结的纳入规则处理全体被试；样本内状态分析与时间留出预测分别保存参数和评价边界。
\item \term{最后生成状态、口述与论文产物。}固定参数与种子来源，完成状态精度和恢复检查，再作规则解释与机制比较。
\end{enumerate}

正式并行使用当前可用的128进程预算，不主动预留核；每个计算进程一个数值线程，只有一层进程并行。
实际进程数不超过当前就绪任务数。并行不能消除逐试次递推和搜索决策之间的依赖，也不应为了占满CPU而增加无必要的粒子或候选。

每个被试至少保存：数据和源码指纹、引擎配置、候选来源、阶段种子和预算、评分掩码、逐轮改进、独立复核、最终参数与近优候选、输出精度、停止原因和未解决问题。
现有结果只读保留，新试验写新目录；只有输入、源码和环境核对通过才能续用数值缓存。

\noindent\term{本稿的主要入口}
\begin{itemize}\setlength{\itemsep}{0pt}\small
\item \href{model_0826.tex}{模型定义与条件适配}；\href{../../workflows/README.md}{工作流与复现命令}。
\item 结果报告：\href{../../../../results/model_0826/simplified_fit_pilot_20260917/README.md}{简化拟合试点}；\href{../../../../results/model_0826/adaptive_effort_pilot_20260917/README.md}{自适应计算试点}；\href{../../../../results/model_0826/search_stopping_pilot_20260918/README.md}{本轮搜索与停止检查}（含配置、数值与验证）。
\item \href{../superpowers/specs/2026-09-01-model-0826-recovery-design.md}{2026-09-01 恢复设计说明}；\href{../../../../CategoryLearning_codes/figures/fig2/ORAL_ALIGNMENT_AUDIT.md}{Fig2 口述编码审查}。
\end{itemize}
\end{document}
